\section{Discussion}
\label{sec:discussion}

The accuracy-spike curve in \S\ref{sec:exp-main} does not separate the proposed
regulariser from Plain $L_2$: both arms land within run-to-run noise of each other
on that curve. \S\ref{sec:exp-analysis} shows the curve is nevertheless not the
whole comparison -- at matched spike counts, active-neuron usage separates the
two arms, 13.1\% fewer for the proposed regulariser. A practice of
comparing spike regularisers by their accuracy-spike curve alone would miss this
gap entirely.

The regulariser was designed around a hypothesis of lateral inhibition -- that
weighting the spike-count penalty by per-neuron activity would make firing
neurons actively suppress their neighbours, reproducing a winner-take-all
dynamic without changing the network architecture. That did not happen. We
measure the correlation between one neuron's firing and its neighbours'
suppression at $+0.61$ to $+0.69$, which is not distinguishable from a
uniform-weighting control at $+0.642$. The concentration reported in
\S\ref{sec:exp-analysis} is real and measured, but it is concentration of
firing onto fewer neurons, not evidence of the suppression mechanism the
method was designed to produce.

The per-neuron term that the regulariser adds to each active neuron's
weighting is a designed, monotone update rule, not the derivative of the loss
stated in \S\ref{sec:method}. The two agree in intent -- both push weighting
away from neurons that are already firing -- but not in form. None of the
results in \S\ref{sec:exp-main}--\S\ref{sec:exp-analysis} depend on treating
the update as a literal gradient; we note the distinction here because a
reader checking the method against its stated objective would otherwise
expect the two to match, and we leave the derivation itself to
\S\ref{sec:method}.

% GATE: G2
Judged only over training epochs 51--200, and only against a fixed-$\lambda$
comparison arm, the ratio controller removes 18.7--20.3\% more spikes than
the fixed setting without a matching loss in accuracy. We do not generalise
this beyond the stated window: the fixed-$\lambda$ arm currently has one run
at $\lambda=3\times10^{-7}$ and one at $\lambda=5\times10^{-7}$, too few to
separate the controller's trajectory over training from its time-averaged
strength, which differs from the fixed comparison value. Isolating the two
needs eight further runs, matched by landing point and by average $\lambda$
respectively, not yet run.

A related question is why $\lambda$ moves the way it does. It drops 36.9\% from
epoch 200 to epoch 201, from $4.71\times10^{-7}$ to $2.97\times10^{-7}$,
coinciding with CutMix being turned off; whether it also coincides with a
specific point in the learning-rate schedule has not been checked. If it does, the ratio
controller would be indirectly coupled to the training schedule in the same
way an earlier, abandoned LR-linked variant of this method was, and the
claim that regularisation strength requires no external tuning would need
qualifying. We do not claim schedule-independence until this is checked.

We do not use a SynOps-based energy estimate to claim an advantage over the
comparison arms. Per-spike SynOps is 1,505 for the proposed regulariser
against 1,488 for Plain $L_2$, a 1.1\% difference this sample cannot resolve,
and the two arms distribute spikes across layers too similarly to build an
energy table that separates them. The defensible statement is that, at the
same SynOps count, the proposed regulariser uses fewer neurons -- not that it
uses less energy.

The two codebases behind this paper's results use different surrogate
gradients. The TensorFlow codebase (ResNet-19, \S\ref{sec:exp-main}--%
\S\ref{sec:exp-analysis}) uses an asymmetric surrogate with bias 0.8 and
window width $\alpha=0.5$ -- the width is the shared library default and is
not overridden by the training configuration these runs use -- while the
QKFormer codebase uses SpikingJelly's \texttt{Sigmoid()}. Porting the
asymmetric surrogate to the QKFormer codebase has been verified against the
TensorFlow implementation (maximum absolute difference $0.00\times10^{0}$),
but it is not enabled there by default, since turning it on would discard
fifteen existing QKFormer runs. Each codebase applies its own choice
uniformly across every arm, so within-codebase comparisons remain valid; we
record the difference because it is inherited from each repository, not a
contribution of this work.

During development we found a porting bug in the QKFormer regulariser: an
axis mismatch caused two of 35 layers to be regularised 2.75$\times$ more
strongly than intended, a 5.9\% difference in the overall objective. We
corrected it and re-measured.

Outside \S\ref{sec:exp-analysis}, this paper's tables and reported figures
are built entirely from eval/test-mode spike counts, never training-mode
logging. We standardise on eval/test throughout because the two protocols
disagree by a wide enough margin in the QKFormer codebase
(\S\ref{sec:exp-setup}) that leaving the choice implicit would risk comparing
two arms under different logging modes without a reader noticing; the one
exception, \S\ref{sec:exp-analysis}'s
active-neuron regression, is marked as training-mode throughout that section
because the per-neuron logging it needs does not exist under eval/test.

``Silent neuron'' means two different things in this paper, and only one of
them is measured. \S\ref{sec:exp-analysis} reports the fraction of neurons
quiet on a given 100-image batch, which underlies every active-neuron and
concentration figure here. The framing that the regulariser turns off
unnecessary neurons, however, implicitly invokes the stronger sense --
neurons that never fire across the entire dataset, which would support a
structural pruning argument. We have not measured that: doing so needs a
full 10,000-image test-set pass (gate G5), and until it exists we do not
equate the two.

\subsection{Limitations}
\label{sec:limitations}

\begin{itemize}\itemsep0pt
  \item Every condition in \S\ref{sec:exp-main} has $n=4$ runs, and no single
  condition-pair comparison reaches significance on its own. The case for
  using fewer active neurons rests on the consistency of its direction across
  the regression in \S\ref{sec:exp-analysis} ($n=21$), not on any one
  pairwise test.
  \item The comparison in \S\ref{sec:exp-main} does not establish that the
  proposed regulariser is more accurate than Plain $L_2$: the
  regression-fitted offset between the two is within run-to-run noise
  ($t=1.41$); reaching conventional significance at this effect size would
  need approximately 48--60 runs, roughly 13 days of compute at the run
  lengths used here, which we have not run and do not plan to.
  \item The neuron-usage concentration reported in \S\ref{sec:exp-analysis}
  is confined to early and middle layers; from conv3\_block2 onward the
  compared arms are not distinguishable, so the claim does not extend to the
  network as a whole.
  \item We have not validated any of these results on neuromorphic hardware.
  Gate G7 (Speck measurement) is unmet, so this paper reports simulation
  only, and the neuron-usage gain in \S\ref{sec:exp-analysis} has not been
  shown to translate into a measured hardware benefit.
  \item Four of the training configuration files that \texttt{run\_paper.py}
  points to sit outside version control, excluded by a \texttt{.gitignore}
  rule for underscore-prefixed directories. Reproducing the exact runs
  behind this paper's tables currently requires access to those files,
  which is unresolved as of this writing.
\end{itemize}
